% Generated draft: requires mathematical and source review.
\documentclass[11pt]{article}
\usepackage[margin=25mm]{geometry}
\usepackage{fontspec}
\setmainfont{Latin Modern Roman}
\usepackage{amsmath,amssymb,mathtools,unicode-math}
\setmathfont{Latin Modern Math}
\usepackage{enumitem,needspace,xurl,hyperref}
\setlength{\parindent}{0pt}
\setlength{\parskip}{5pt}
\setlength{\emergencystretch}{3em}
\newcommand{\N}{\mathbb{N}}
\newcommand{\Z}{\mathbb{Z}}
\newcommand{\Q}{\mathbb{Q}}
\newcommand{\R}{\mathbb{R}}
\newcommand{\C}{\mathbb{C}}
\newcommand{\F}{\mathbb{F}}
\newcommand{\A}{\mathbb{A}}
\newcommand{\PP}{\mathbb{P}}
\newcommand{\E}{\mathbb{E}}
\newcommand{\eps}{\varepsilon}
\newcommand{\dd}{\,\mathrm{d}}
\begin{document}
\section*{UnsolvedMath ID: 30003203}
\textbf{Source problem code:} catalogue entry
\subsection{Definitions and mathematical statement}
Write \(M(\mu,m)=\Gamma(\mu,m)''\) for the free Araki--Woods factors referred to in the supplied record. In the usual spectral description, \(\mu\) describes the spectral measure class of the underlying orthogonal representation of \(\R\), and \(m\) describes its spectral multiplicity. The record does not specify the admissible measures, multiplicity functions, dimension restrictions, or separability assumptions; these parameter conventions require confirmation before a fully specified parameterwise classification can be stated.

For terminology, let \(U:\R\to\mathcal O(H_{\R})\) be a strongly continuous orthogonal representation on a real Hilbert space, and let \(U^{\C}\) be its complexification on \(H_{\C}\). The representation is almost periodic if
\[
H_{\C}=\overline{\operatorname{span}}\bigl\{\xi\in H_{\C}:\text{ there exists }\lambda\in\R\text{ such that }U_t^{\C}\xi=e^{it\lambda}\xi\text{ for every }t\in\R\bigr\}.
\]
Here almost periodicity refers to the generating representation. Its failure allows both continuous spectral data and mixtures of atomic and continuous spectral data.

\textbf{Classification problem.} Classify the free Araki--Woods factors \(M(\mu,m)\) beyond the almost-periodic case. Under the conventional interpretation of classification as classification of von Neumann algebras up to normal unital \(*\)-isomorphism, seek an explicitly described invariant \(I\) such that, for every two admissible parameter pairs \((\mu,m)\) and \((\nu,n)\) in the intended classification domain,
\[
M(\mu,m)\cong M(\nu,n)
\quad\Longleftrightarrow\quad
I(\mu,m)=I(\nu,n).
\]
The intended domain must include generating representations that are not almost periodic. The record does not specify whether the requested formulation must also describe every comparison between almost-periodic and non-almost-periodic generating data. It also does not explicitly specify the isomorphism convention; the convention above does not require preservation of a distinguished state and requires source confirmation.

No candidate complete invariant or conjectural equivalence criterion is supplied. The original formulation explicitly emphasizes the absence of even a conjectural classification statement. Thus the displayed equivalence specifies the requested form of a classification, rather than asserting a proposed theorem.
\subsection{Short English statement}
Find a complete classification of free Araki--Woods factors extending beyond those generated by almost-periodic representations. The goal is an explicit invariant that agrees exactly when the factors are isomorphic. The supplied record gives neither a candidate invariant nor a complete specification of the parameter domain and equivalence convention. Its partially solved status and claimed absence of a general classification are supplied dated assessments requiring independent review.
\subsection{Sources}
\begin{itemize}
\item Supplied local UnsolvedMath registry record, permanent ID \(30003203\), title: Classification of General Free Araki--Woods Factors; exact source problem code: \texttt{catalogue entry}. Collection: Oberwolfach Reports Open Problems. The record lists proposed year 2016 and supplies no individual proposer.
\item Supplied report citation: \(C^*\)-Algebras (2017), Oberwolfach Reports, DOI 10.4171/owr/2016/40, pp.\ 2269--2345. \url{https://doi.org/10.4171/owr/2016/40}. The citation's 2017 label and the record's 2016 proposed year are retained without resolving their relationship.
\item Supplied literature source, labelled ``Houdayer Shlyakhtenko and Vaes Free Araki Woods factors'': \url{https://arxiv.org/abs/1605.06057}.
\item Supplied literature source, labelled ``OWR Operator Algebras report'': \url{https://ems.press/content/serial-article-files/46645?nt=1}.
\item Literature assessment dated exactly \texttt{2026-08-21}: the supplied curation labels the problem partially solved and reports that a large family of non-almost-periodic free Araki--Woods factors has been completely classified, while no general classification beyond the almost-periodic case is known. This is supplied dated curation, not an independently verified literature assessment.
\item The supplied statement status is \texttt{corrected\_verified}. Its verification field reports that the official report was checked, that classification beyond the almost-periodic case was identified as open, and that no conjectural complete classification was available. These are claims of the supplied record; no source visit or independent verification is asserted here.
\end{itemize}
\end{document}
